\paragraph{Results.} The alternative-workflow limits were 75.18\%, 53.65\% and
27.99\% against 75.13\%, 53.65\% and 27.99\% in the original workflow on the same
backgrounds (Table~\ref{tab:coef-primary}). The leading-order limit changed by
less than 0.3 percentage points across all 22 rows; 21 of the 22 differences
were smaller than 0.06 percentage points, and the largest, $+0.27$ points, was
at .95$\to$.99, $\rho=0$, $e=1$, where $K_1$ is smallest
(Table~\ref{tab:surface-diff}; Fig.~\ref{fig:surface-diff}). All three
prespecified predictions held (Table~\ref{tab:predictions}). Site clearing
cannot act at $p=0$, but the isolation rule changed every background's history
and lowered every coefficient (see below), so ratio invariance was not
guaranteed; prediction 1 assumed that the rule would move $K_0$ and $K_1$
nearly in proportion. The predictions were therefore genuine but low-risk, with
generous tolerances (2 percentage points against a realized 0.05, and .05
against .005 for the protocol's two-point crossing statistic); the measured size of the change is the result. Several paired intervals exclude zero; with differences of 0.001
to 0.3 percentage points they reflect the strong within-background pairing, and
they carry no substantive interpretation.

